\section{Backward demand propagation and finite-use composition}
\label{sec:demand}

Precision and resource requirements flow opposite to data.  A consumer knows
how accurately it must answer; from that it computes how accurately each
predecessor must be observed.  The graph then runs forward at the resulting
demands.  This reverse--forward architecture turns the phrase ``choose all
approximations sufficiently fine'' into a primitive-recursive program with a
proof obligation.

\subsection{Static demand graphs}

Let the graph nodes be \(0,\ldots,N-1\) in topological order.  For every edge
\(i\to j\), the consumer \(j\) has a monotone demand transformer
\[
 \delta_{j,i}:\N\times Z\longrightarrow\N,
\]
where \(Z\) denotes fixed input parameters and finite summaries known before
the producer runs.  Let \(q_i\) be the externally requested demand at node
\(i\), usually zero except at exported outputs.  Define demands in reverse
order by
\begin{equation}
 d_i=\max\left(q_i,
       \max_{i\to j}\delta_{j,i}(d_j,z)\right).
 \label{eq:backward-demand}
\end{equation}
The empty maximum is zero.

Because \(N\) and the adjacency lists are fixed, this is a finite composition
of PR functions and maximum, hence primitive recursive.  One may compile it
as straight-line syntax or as a fixed-length reverse table fold.

\begin{theorem}[Soundness and leastness of backward demands]
\label{thm:demand-least}
Assume every \(\delta_{j,i}\) is monotone in its demand argument.  The vector
\(d\) defined by \eqref{eq:backward-demand} satisfies:
\begin{enumerate}
\item \(d_i\geq q_i\) for every node;
\item \(d_i\geq\delta_{j,i}(d_j,z)\) for every edge \(i\to j\);
\item if \(e\) satisfies the preceding two families of inequalities, then
\(e_i\geq d_i\) for every \(i\).
\end{enumerate}
Thus \(d\) is the pointwise least safe demand assignment.
\end{theorem}

\begin{proof}
The first two statements follow directly from the defining maxima.  For
leastness, use reverse induction on \(i\).  Suppose \(e_j\geq d_j\) for all
successors \(j>i\).  Monotonicity gives
\(\delta_{j,i}(e_j,z)\geq\delta_{j,i}(d_j,z)\).  Since \(e\) is safe,
\(e_i\) dominates \(q_i\) and every left-hand edge demand evaluated at
\(e_j\), hence it dominates the maximum defining \(d_i\).
\end{proof}

For formalization, the reverse induction is implemented as a forward primitive
recursion on the reverse index \(r=N-1-i\).  A table invariant states that all
already-filled rows equal their defining maxima.  The relevant packet already
contains checked object derivations for the zero and successor equations of a
backward-demand compiler; these should be reused rather than reimplemented.

\subsection{Local soundness under computed demands}

A deterministic node contract is parameterized by a requested output demand:
\begin{equation}
 \left(\bigwedge_{i\to j}
  \Valid_i(x_i,\delta_{j,i}(q,z))\right)
 \wedge \Cert_j(\vec x,q,c)
 \longrightarrow
 \Valid_j(E_j(\vec x,q,c),q).
 \label{eq:node-contract}
\end{equation}
If predecessor \(i\) is produced validly at \(d_i\), restriction coherence and
Theorem~\ref{thm:demand-least} give validity at
\(\delta_{j,i}(d_j,z)\).  Instantiating \eqref{eq:node-contract} at
\(q=d_j\) therefore proves node \(j\)'s output valid at \(d_j\).

\begin{theorem}[Forward graph validity]
\label{thm:forward-validity}
Suppose every source node is valid at its computed demand, every node has a
local contract of the form \eqref{eq:node-contract}, and restriction is sound.
Then a topological forward run produces every node valid at its computed
demand.  In particular every exported node is valid at its external request.
\end{theorem}

\begin{proof}
Induct in topological order.  At node \(j\), all predecessors are valid at
their computed demands by the induction hypothesis.  Restrict them to the
local demands, use the node certificate, and apply the local contract.  The
external conclusion follows from \(d_j\geq q_j\) and restriction.
\end{proof}

This theorem is intentionally modular: the backward pass establishes only
numerical inequalities; the forward proof establishes only validity.  The two
are joined by restriction coherence.

\subsection{Multiple resources}

Demand need not be a single precision.  A resource vector can contain:
\[
 (\text{precision},\text{degree},\text{coefficient height},
   \text{matrix dimension},\text{iteration fuel},\text{exception budget}).
\]
Use the componentwise preorder and componentwise join.  A demand transformer
returns a resource vector, and the reverse equation uses joins instead of
maximum.  The leastness proof is unchanged for any finite product of
join-semilattices with monotone transformers.  Since the number of coordinates
is fixed, all joins are PR.

This generalization is useful in algebraic geometry.  A consumer may need a
jet of order \(m\), coefficients at height \(H\), and a denominator apartness
bound \(a\).  Flattening these into one enormous scalar obscures the proof;
keeping a fixed tuple makes the contract readable while remaining a natural
number code internally.

\subsection{Value-dependent demands}

Some algorithms appear to demand precision based on a value produced later.
Three constructively legitimate repairs cover most cases.

\paragraph{Safe envelope.}
Replace the exact value-dependent demand by a monotone upper bound computed
from input parameters and already available magnitude certificates.  This is
the preferred method for arithmetic operations.

\paragraph{Challenge summary.}
Run a coarse finite pass that computes a bounded summary \(z\), such as an
estimated magnitude class, a pivot location, or a candidate degree.  The
backward demand pass is then allowed to depend on \(z\), and a second forward
pass replays the computation at the resulting demands.  The summary's
correctness is itself a trace-node contract.

\paragraph{Relational retention.}
If no stable summary can be selected, carry every finite candidate and let the
consumer quantify over the cloud.  Backward viability can eliminate candidates
that cannot extend to the requested horizon.  No value-dependent global
selector is needed.

What is not legitimate is a circular definition in which the producer needs a
demand that depends on an exact output obtainable only after running the
producer at that demand.  A PCOG must break the circle by one of the three
methods above.

\subsection{Two-pass trace invariant}

A fixed graph can be executed by one state machine with two phases:
\begin{enumerate}
\item reverse phase: fill the demand table;
\item forward phase: fill the value and certificate table.
\end{enumerate}
A state encodes the phase, program counter, demand table, output table, and
certificate flags.  Its quantifier-free invariant says:

\begin{itemize}
\item the program counter lies in the phase's fixed range;
\item every completed reverse row satisfies \eqref{eq:backward-demand};
\item every completed forward row is valid at its recorded demand;
\item untouched rows contain fixed defaults;
\item the demand table is unchanged during the forward phase;
\item every certificate flag corresponds to a checker returning one.
\end{itemize}

All universal statements over completed rows are represented by bounded
Boolean folds.  The step proof dispatches on the fixed program-counter range.
Theorem~\ref{thm:trace-reflection} then proves the invariant at the terminal
state.

For maintainability, the recommended implementation uses separate trace nodes
for reverse and forward passes, connected by the static proof graph.  The
single-state-machine presentation is nevertheless useful as a completeness
argument: the entire finite graph can, in principle, be reduced to one
primitive recursion and one quantifier-free induction.

\subsection{Composition law for contracts}

Suppose \(F:X\to Y\) has demand \(\delta_F\) and \(G:Y\to Z\) has demand
\(\delta_G\).  The composite demand is
\[
 \delta_{G\circ F}(p)=\delta_F(\delta_G(p)).
\]
If each local validity theorem is embedded, the composite theorem is obtained
by instantiating \(F\)'s contract at \(\delta_G(p)\), then applying \(G\)'s
contract at \(p\).  This is ordinary cut.

For a DAG, reverse propagation computes the closure of this composition across
all paths and takes the maximum where a producer is shared.  Therefore a
shared node is evaluated once at the greatest demand of all consumers.  Its
result can be restricted independently along each outgoing edge.

\subsection{Precision accounting as proof data}

The demand vector should be retained in certificates and theorem statements,
not erased after execution.  It supports three audits:

\begin{enumerate}
\item \emph{dependency audit}: which input observations were actually read;
\item \emph{complexity audit}: what bounded fuel and table sizes were used;
\item \emph{soundness audit}: whether every local theorem was instantiated at
a demand no larger than the data supplied.
\end{enumerate}

This is especially useful for large theorem packets.  A machine-readable graph
certificate can list each node, its demand expression, local theorem name, and
predecessors.  The final Lean proof term remains authoritative, but the graph
certificate makes the construction reviewable.
